%%=============================================================================
%% Samenvatting
%%=============================================================================

% TODO: De "abstract" of samenvatting is een kernachtige (~ 1 blz. voor een
% thesis) synthese van het document.
%
% Een goede abstract biedt een kernachtig antwoord op volgende vragen:
%
% 1. Waarover gaat de bachelorproef?
% 2. Waarom heb je er over geschreven?
% 3. Hoe heb je het onderzoek uitgevoerd?
% 4. Wat waren de resultaten? Wat blijkt uit je onderzoek?
% 5. Wat betekenen je resultaten? Wat is de relevantie voor het werkveld?
%
% Daarom bestaat een abstract uit volgende componenten:
%
% - inleiding + kaderen thema
% - probleemstelling
% - (centrale) onderzoeksvraag
% - onderzoeksdoelstelling
% - methodologie
% - resultaten (beperk tot de belangrijkste, relevant voor de onderzoeksvraag)
% - conclusies, aanbevelingen, beperkingen
%
% LET OP! Een samenvatting is GEEN voorwoord!

%%---------- Nederlandse samenvatting -----------------------------------------
%
% TODO: Als je je bachelorproef in het Engels schrijft, moet je eerst een
% Nederlandse samenvatting invoegen. Haal daarvoor onderstaande code uit
% commentaar.
% Wie zijn bachelorproef in het Nederlands schrijft, kan dit negeren, de inhoud
% wordt niet in het document ingevoegd.

\IfLanguageName{english}{%
\selectlanguage{dutch}
\chapter*{Samenvatting}
\lipsum[1-4]
\selectlanguage{english}
}{}

%%---------- Samenvatting -----------------------------------------------------
% De samenvatting in de hoofdtaal van het document

\chapter*{\IfLanguageName{dutch}{Samenvatting}{Abstract}}

Deze bachelorproef onderzoekt de sim-to-real prestaties van een autonome magazijnrobot die volledig virtueel is getraind met NVIDIA Isaac Sim en ROS 2. Het centrale vraagstuk is hoe goed het gedrag en de prestaties van een robot, ontwikkeld in een digitale omgeving, overeenkomen met zijn functioneren in een echte magazijncontext.

De robot wordt vanaf nul ontworpen en gebouwd, inclusief hardware (chassis, aandrijving, hefmechanisme) en softwaremodules (ROS 2-nodes voor communicatie en besturing). In de simulatieomgeving worden navigatie, mapping en obstakelvermijding getraind en getest, waarna dezelfde taken fysiek worden uitgevoerd in een gecontroleerde magazijnsetting. De prestaties worden kwantitatief vergeleken op metrics zoals taakduur, afgelegde afstand, botsingen en kaartnauwkeurigheid.

Verwacht wordt dat er een duidelijke kloof ontstaat tussen simulatie en realiteit, veroorzaakt door sensorruis, variabele lichtomstandigheden en moeilijk te simuleren factoren zoals vloerfrictie. Toch biedt de simulatie-first aanpak aanzienlijke voordelen: kortere ontwikkeltijd, lagere kosten en snellere iteraties. De meerwaarde voor IT-professionals en bedrijven ligt in praktische richtlijnen voor virtuele training, inzicht in beperkingen van sim-to-real overdracht en aanbevelingen voor het optimaliseren van magazijnautomatisering.